\documentclass[11pt]{article}
\usepackage[margin=0.83in]{geometry}
\usepackage{amsmath,amssymb,booktabs,array,microtype}
\usepackage[T1]{fontenc}
\usepackage{lmodern}
\usepackage[hidelinks]{hyperref}
\usepackage{tikz}
\usetikzlibrary{arrows.meta}
\setlength{\parindent}{0pt}
\setlength{\parskip}{6pt}
\newcommand{\avg}[1]{\overline{#1}}
\begin{document}
\begin{center}
{\Large Housing, borrowing, and family formation}\\[4pt]
{\large A proposed repair of the illustrative model}\\[4pt]
September 9, 2026
\end{center}

The most promising change is to distinguish \emph{cash at purchase} from
\emph{earnings received during the period}. This permits a simple equilibrium
in which young households rent small homes, save from their earnings, and buy
larger homes when old. A planner choosing consumption and housing gives more
housing to the young. With a further explicit restriction on child goods
costs, the planner also chooses higher fertility.

This is a proposal for discussion. It preserves the two ages, logarithmic
utility, child goods and space costs, the two upper housing limits, property
taxes, and income--wealth heterogeneity. It changes the financial calendar,
allows tenure to change with age, and omits ownership tastes in the first
illustration. The main seminar deck has not been changed.

\textbf{1. Environment.}
Households are young, then old. Young households choose completed fertility
$n$. Their preferences remain
\[
u^y=\log(c^y-\chi n)+\alpha\log(h^y-\kappa n)+\vartheta\log n,
\qquad
u^o=\log c^o+\gamma\log h^o+\omega_B\log e.
\]
Lifetime utility is $u^y+\beta u^o$. The estate $e$ includes the house's sale
value at death. For the allocation comparison, set $\alpha=\gamma$.
All preference and child-cost coefficients are positive.

Rental homes have size $h\le H_R$ and owned homes have size $h\le H_O$,
where $H_R<H_O$. There is a common divisible stock $\bar H$. This retains
the original common-stock interpretation: dwelling sizes and ownership
can change, subject to these menus. It is not a fixed inventory of
nonconvertible rental and owner units.

\textbf{Cash and earnings.}
A young household has cash $w_i$ at purchase, including every receipt already
received. Its later working income is $Y_i$. These may differ across
households. A mortgage finances at most the fraction $\phi$ of the purchase.
It is repaid from working income before the household enters old age;
financial assets at that boundary must be nonnegative. Retirement labor
income is zero in the example below. Old households can sell their homes and
choose either tenure.

Both current ages consume at the end of the same period. Old households
leave their estates at that date. A young household's own old consumption
comes one period later. Thus the dated planner compares goods and housing
delivered on a common calendar.

\textbf{Why this matters.}
A household can afford housing payments out of its earnings and still lack
the deposit at purchase. Putting an entire long period's income in the bank
before the purchase weakens that distinction. The proposal changes when
income arrives; it does not exclude income that has already been received.

\newpage
\textbf{2. Household budgets.}
Let $q\in(0,1)$ be the price of a bond paying one unit at period end, $P$ the
stationary house price, and $\tau>0$ the property-tax rate. Let $R=1/q$ denote
the gross bond return, and write $D=1-q+q\tau$. End-of-period rent per unit is
$DP/q$.
This is a small open economy with outside-sector bond and intermediary finance;
warm-glow estates do not determine the next entrant's cash.
If each age has mass $N$, the common upfront tax rebate is
\[
T=\frac{q\tau P\bar H}{2N}.
\]
The government finances this rebate against the period-end property-tax
receipts. It introduces no additional goods.

A young owner chooses consumption, housing and next-period financial wealth
$a'$ subject to
\[
\boxed{\begin{aligned}
qc^y+qa'+(1+q\tau)Ph^y&=w_i+T+qY_i,\\
(1-\phi)Ph^y&\le w_i+T,\qquad a'\ge0.
\end{aligned}}
\]
The second line is the deposit requirement. In particular, $\phi=.8$ means
that cash must cover 20\% of the purchase. A young renter instead has
\[
qc^y+qa'+DP h^y=w_i+T+qY_i,
\qquad a'\ge0,\quad h^y\le H_R.
\]
This timing requires the home to be chosen before the later earnings arrive.
A purchase after that receipt would be a subsequent housing decision.

An old household's initial resources include its financial wealth, any
inherited title, and the rebate:
\[
Z=a'+P h^y\mathbf 1\{\text{previously owner}\}+T.
\]
An old owner has budget and estate
\[
qc^o+qe+DP h^o=Z,\qquad
e=a^e+Ph^o,\qquad a^e\ge0.
\]
Here $a^e$ is the financial estate delivered at death. An old renter has the
same consolidated budget, with $e=a^e$ and $h^o\le H_R$.
The owner's financial lower bound implies $e\ge Ph^o$. Old households need
not leave liquid assets in the equilibrium used below.

\textbf{The deposit can exclude large homes without an owner minimum.}
If
\[
\frac{w_i+T}{(1-\phi)P}<H_R,
\]
every home the young household can buy is smaller than the largest rental.
A renter can replicate such an owner's consumption and future total wealth
by saving the house's value in bonds. With mobile future tenure and no
ownership taste, buying that smaller home has no advantage. If rental demand
reaches $H_R$, the young household rents there and would like more space.
This is the joint role of the deposit and rental-size constraints.

\newpage
\textbf{3. An equilibrium with the required allocation.}
First hold fertility fixed, with $0<\kappa n_i<H_R$ and
$\nu\avg n=1$. Normalize each age mass to one and write
$r=H_R$, $H=\bar H/N$, and $m=H-r$. The proposed equilibrium has all young
households renting $r$ and mean old housing $m>r$.

Define the old utility weight and housing coefficient by
\[
K=1+\gamma+\omega_B,\qquad
a=\min\left\{\frac{\gamma}{D},
                  \frac{\gamma+\omega_B}{1+q\tau}\right\}.
\]
The two terms cover positive financial estates and a binding estate floor.
In either case, an uncapped old owner has
\[
c_i^o=\frac{Z_i}{qK},\qquad h_i^o=\frac{a Z_i}{KP}.
\]
For young net resources $b_i=w_i/q+Y_i-\chi n_i$, set $G=1+\beta K$.
The equilibrium price and rebate are explicit:
\[
\boxed{\begin{gathered}
J=Gm+a\beta\left[\frac Dq r-\frac{(1+q)\tau H}{2}\right],\\
P=\frac{a\beta\avg b}{J},\qquad T=\frac{q\tau PH}{2}.
\end{gathered}}
\]
The individual allocation is then
\[
\boxed{\begin{aligned}
x_i\equiv c_i^y-\chi n_i
 &=\frac{b_i+(1+q)T/q-(DP/q)r}{G},\\
a_i'&=\beta Kx_i-T,\qquad
c_i^o=\frac\beta qx_i,\qquad
h_i^o=\frac{a\beta x_i}{P}.
\end{aligned}}
\]
The next page states every inequality needed to verify this equilibrium.
They are functions of primitives because $P$ and $T$ are already given
by the displayed formulas.

\textbf{A finite region where all the inequalities hold.}
One analytical family has
\[
\begin{gathered}
q=.5,\quad \phi=.8,\quad r=1,\quad H=2.5,\quad H_O=2,\\
\alpha=\gamma=1,\quad\omega_B=.25,\quad
.3\le\beta\le1,\quad .01\le\tau\le.02,\\
\left|b_i/\avg b-1\right|\le .05,
\qquad \avg b\ge45\max_i w_i.
\end{gathered}
\]
Every young household rents one unit and saves; each old household owns
between $1.405$ and $1.595$ units. The old financial estate is zero.
These conclusions hold throughout the stated region, by analytical
inequalities. The numbers do not stand in for a numerical equilibrium proof.

The substantive restriction is low cash relative to later stage income.
This is a sufficient example, with limited resource dispersion and no tenure
tastes, not an empirical assessment of the parameter region.

\newpage
\textbf{4. Complete verification of the equilibrium statement.}
Let $b_e=\max\{\omega_B,qa\}$ and $G_R=1+\beta(1+\omega_B)$.
For $t>1$, define the utility advantage of the owner continuation by
\[
\Delta(t)=\beta\gamma\log t
 +\beta\omega_B\log(b_e/\omega_B)
 -G_R\log\left(\frac{G-a\beta D/t}{G_R}\right).
\]
This compares the optimized lifetime plan ending as an owner with the
optimized plan renting in both ages; it is not a marginal-utility assumption.

\textbf{Proposition.}
Suppose $J>0$, $\avg b>0$, and, at the explicit $P,T,x_i,h_i^o$ above,
every type satisfies
\[
\boxed{\begin{aligned}
&\frac{w_i+T}{(1-\phi)P}<r,
&&\text{purchase cash limits the young;}\\
&x_i>0,\quad \beta Kx_i>T,
&&\text{consumption and saving are positive;}\\
&\alpha x_i>\frac{DP}{q}(r-\kappa n_i),
&&\text{the young rental limit binds;}\\
&r<h_i^o<H_O,
&&\text{the old owner choice is feasible;}\\
&\Delta(h_i^o/r)>0,
&&\text{owning when old is preferred.}
\end{aligned}}
\]
Then the displayed prices and allocations form a stationary competitive
equilibrium. All young rent at their size limit, save, and buy larger homes
when old. Prices are unique within this regime. Other equilibrium regimes
have not been excluded.

\textbf{Proof.}
With a future owner, the old value is $K\log Z$ plus a constant. The young
problem is strictly concave; its saving condition gives $Z_i=\beta Kx_i$.
The stated cap and feasibility inequalities verify its optimum. The old
estate floor is included in that value, even when it binds.

The alternative lifetime renter problem is jointly concave. Both rental
limits bind under the displayed restrictions, and its adult consumption is
\[
x_i^{RR}=\frac{Gx_i-DPr}{G_R}.
\]
Subtracting its maximized value gives $\Delta(h_i^o/r)$. The deposit
replication on page 2 rules out a better plan starting as an owner. Finally,
the price formula gives $\avg h^o=m$, and $2T=q\tau PH$ balances the fiscal
account. Rental intermediaries earn the bond return and all titles are
accounted for.

\textbf{Chosen fertility.}
At the rental limit, a private choice of fertility solves
\[
\frac{\vartheta}{n_i}
=\frac{\chi}{x_i}+\frac{\alpha\kappa}{r-\kappa n_i}.
\]
Substitution for $x_i$ gives a quadratic. The supporting calculation extends
the displayed family by imposing small child goods costs and verifying
tenure preference over \emph{all} feasible fertility choices. A one-dimensional
monotone equation then imposes replacement fertility. Heterogeneous fertility
need not be replaced by a representative household.

\newpage
\textbf{5. The full dated planner.}
The planner reallocates consumption and housing among the households alive
at the start of a period. First hold their fertility fixed. It chooses
current tenure as well as home size, respects the housing stock and the
size menus, and can relax private financial constraints.
Individual young continuation wealth and old estates stay fixed. Future prices and
opportunity sets are held fixed in this dated comparison. Thus their future utility
terms are constant in the planner's problem.

Let $\mathcal C$ be total goods delivered to the two current ages, per unit
of young-cohort mass. With $\alpha=\gamma$, its problem is
\[
\begin{aligned}
\max_{c^y,c^o,h^y,h^o,\,\text{tenure}}\quad
 &\int_y\!\left[\log(c_i^y-\chi n_i)
                  +\alpha\log(h_i^y-\kappa n_i)\right]
 +\int_o\!\left[\log c_j^o+\alpha\log h_j^o\right]\\
\text{subject to}\quad
 &\int_y c_i^y+\int_o c_j^o=\mathcal C,\qquad
 \int_y h_i^y+\int_o h_j^o=H,
\end{aligned}
\]
with positive adult consumption and the stated size menus. Since ownership
has no direct taste cost, the planner can assign any household a feasible
owned home up to $H_O$. This permission to change tenure is essential: a
young household whose rental tenure were fixed could not exceed $H_R$.

The solution equalizes adult consumption. With slack planner housing caps,
\[
X=\frac{\avg x^y+\avg c^o}{2},\qquad
s=\frac{H-\kappa\avg n}{2},\qquad
h_i^{y,F}=s+\kappa n_i,\quad h_j^{o,F}=s.
\]
With binding caps, replace these housing choices by
$\min\{H_O,s+\kappa n_i\}$ and $\min\{H_O,s\}$, and choose $s$ to clear
the stock. Young mean housing is at least $H/2$.

\textbf{Allocation result.}
In the equilibrium on pages 3--4, $h_i^{y,E}=r$ and $H=r+m>2r$.
Consequently,
\[
\boxed{\avg h^{y,F}\ge\frac{r+m}{2}>r=\avg h^{y,E}.}
\]
The full planner allocates more housing to the young. With common fertility
and slack caps, the increase for each young household is
\[
h^{y,F}-h^{y,E}=\frac{m-r+\kappa n}{2}>0.
\]
Under heterogeneous fertility the aggregate claim still holds; an individual
gain is checked from the displayed allocation, not inferred from its mean.

The comparison is financed by current transfers
$q\Delta c_i+DP\Delta h_i$. They sum to zero. Changes in title holdings are
offset by financial claims, including those of rental intermediaries, keeping
young continuation wealth and old net estates unchanged. This is a dated
utilitarian result, not a comparison that omits an initial old cohort.

\newpage
\textbf{6. Fertility: a condition without a restriction on $\beta R$.}
Now let the same planner choose young fertility, valuing current parents'
utility. Children require goods and space in the young period only. Future
parental utility therefore remains fixed under the same continuation-wealth
settlement. No utility for unborn households is added.

Let $n_0$ be mean private fertility and $\avg x$ mean young adult
consumption in an equilibrium where private fertility is optimal. Since
$c_i^o=(\beta/q)x_i$, current goods per young-cohort unit are
\[
\mathcal C=(1+\beta/q)\avg x+\chi n_0.
\]
The joint planner chooses common fertility $n^F$, with adult consumption
and young adult space
\[
X(n)=\frac{\mathcal C-\chi n}{2},\qquad
S(n)=\min\left\{\frac{r+m-\kappa n}{2},\ H_O-\kappa n\right\}.
\]
Its unique interior fertility choice satisfies
\[
\boxed{\quad\frac{\vartheta}{n^F}
 =\frac{\chi}{X(n^F)}+\frac{\alpha\kappa}{S(n^F)}.\quad}
\]
On either housing-cap regime this is a quadratic. Thus consumption and
housing are reoptimized when the planner chooses fertility.

There is an exact test for a rise in mean fertility:
\[
\boxed{\quad n^F>n_0\quad\Longleftrightarrow\quad
\frac{\vartheta}{n_0}>
\frac{2\chi}{(1+\beta/q)\avg x}+\frac{\alpha\kappa}{S(n_0)}.\quad}
\]
If $\beta\ge q$, the equilibrium's old consumption supplies a sufficient
goods cushion and the inequality holds. When $\beta<q$, a simple sufficient
condition, valid under income heterogeneity and binding planner caps, is
\[
\boxed{\quad
\frac{\chi r}{\kappa\avg x}\frac{q-\beta}{q+\beta}
 <\alpha\frac{m-r}{m+r}.
\quad}
\]
The left side measures the goods cost of supporting children when consumption
is equalized across ages. The right side measures the available housing
gain. Cheap child goods or a larger initial housing difference makes a
fertility increase easier. This is not a condition on a credit multiplier.

A stronger condition that works for \emph{every} $\beta>0$ is
\[
\frac{\chi r}{\kappa\avg x}<\alpha\frac{m-r}{m+r}.
\]
In the explicit family, the right side is $.2$. Impose the primitive restrictions
\[
g_i=\frac{w_i}{q}+Y_i,\qquad
\left|\frac{g_i}{\bar g}-1\right|\le .01,\qquad
\frac{\chi r}{\kappa}\le .02\bar g,\qquad
\bar g\ge46\max_i w_i.
\]
These primitive restrictions strengthen the family on page 3. For any specified
$\nu>\kappa/r$, a unique scalar calibration of $\vartheta$ makes mean private
fertility equal $1/\nu$ at the normalized cohort size. The price is then explicit
and individual fertility is quadratic. The proof checks all fertility and tenure
deviations. Throughout this family, $\chi r/(\kappa\bar x)<1/12<1/5$, so the
full joint planner raises mean fertility, including $\beta R<1$. Calibrating
$\vartheta$ to replacement is not a closed-form solution for $\vartheta$ or a
claim of global equilibrium uniqueness.

\newpage
\textbf{7. The transition in the same model.}
Keep the entrant distribution of $(w,Y)$ fixed and exogenous at every date;
there is no endogenous transmission of types through fertility or estates. A
permanent taste shock arrives when the stage opens. With adult household masses
$N_t^y$ and $N_t^o$, demographic accounting is
\[
N^y_{t+1}=\nu\bar n_tN^y_t,
\qquad N^o_{t+1}=N^y_t.
\]
Both positive stationary endpoints have $\bar n^*=1/\nu$, and population means
refer to adult households, so a stationary total is $N^y+N^o=2N$.

In this equilibrium, old households leave only housing in their estates. Write
\[
a=\frac{\gamma+\omega_B}{1+q\tau},\qquad
G=1+\beta K,\qquad
C_0=(1-q)\left(\frac1q+\frac\tau2\right),\qquad
\widetilde G=G-\frac{(1+q)\tau a\beta}{2}.
\]
At replacement fertility, the stationary endpoints satisfy
\[
P^*=\frac{\bar g-\chi/\nu-\widetilde G\bar x^*}{C_0r},
\qquad
N^*=\frac{\bar H}{r+a\beta\bar x^*/P^*}.
\]
Here $\bar x^*$ is the mean from aggregating individual quadratic fertility
choices at replacement. Prices, rebates, and inherited saving are endogenous
along the path.

The local theorem for the explicit endogenous family is as follows. A small
permanent decline in $\vartheta$ lowers mean fertility on impact and lowers the
endpoint $N$. At any finite date on that nearby baseline, a small unexpected
permanent rebated-tax increase raises mean fertility on impact and raises the
endpoint $N$. The path exists uniquely locally in the verified regime and
converges to its nearby stationary endpoint. At the surprise, inherited saving
is held fixed while the actual rebate changes: the current purchase price falls
on impact, whereas the endpoint price rises. Young households continue to rent
at $h=r$; the tax response works through goods, rents, saving, and rebates and
does not implement the planner's full housing upgrade.

Mean inherited financial wealth is $\bar a_t$. Linear old demand in inherited
saving yields an exact three-dimensional aggregate recursion in
$(N_t^y,N_t^o,\bar a_t)$. The aggregate fertility elasticity, defined as the
elasticity of mean fertility with respect to mean adult consumption, satisfies
$\epsilon<.13$, and the root inequalities establish local stability in the
verified regime. The full calculation supporting local convergence is in the
supporting record. Fertility need not be higher at every subsequent date, and
convergence need not be monotone.

\newpage
\textbf{8. What to keep, and what remains to do.}
The proposed cash-flow model gives the clearest combination obtained in this
pass: two ages, solved household choices, a complete stationary example,
the full dated housing result, and a fertility condition with an economic
interpretation. Its substantive choices should be decided before changing
the main deck:

\begin{enumerate}
\item Treat $w$ as cash already received and $Y$ as later working income.
Consumption for both current ages and the old estate are delivered at period
end. This is a change in financial timing, not a change of letters.
\item Let tenure change with age. The first illustration has no ownership
taste; income and wealth still differ across households. The planner may
reassign tenure and works with the original common divisible stock.
\item Repay mortgages before entry into old age, with nonnegative financial
assets at that boundary. This is stronger than permitting a house sale at
that date to finance repayment. The equilibrium constructed here has saving
young renters, rather than a positive mass of realized young mortgages.
\end{enumerate}

The older calculations remain useful alternatives:

\small
\begin{tabular}{p{.24\textwidth}p{.32\textwidth}p{.33\textwidth}}
\toprule
Alternative & What it delivers & Main cost \\
\midrule
Original timing, owner minimum, mobile tenure
& Explicit fixed-fertility price and the intended age allocation
& Its checked family still requires strong old-income and estate demand. \\[6pt]
Three ages, with a mature earning period
& An analytical equilibrium construction and a full planner result; quadratic or cubic household choices
& An additional age and financial restriction; the useful family uses high mature income. \\[6pt]
Original two-age log specification, weaker estate assumptions
& An adult-space welfare result at conventional LTVs
& In the checked moderate-taste family, old gross homes are much smaller than young homes. \\[6pt]
Real moving costs
& A reason for slow resizing after a fertility shock
& The planner also pays those costs; retention alone does not prove misallocation. \\
\bottomrule
\end{tabular}
\normalsize

More generally, concavity alone does not order housing by age. If $v(h,n)$ is
housing utility, child-space increasing differences such as
$v_h(h,n)>v_h(h,0)$ provide the relevant marginal preference ingredient; the log
housing term satisfies it.

\textbf{Remaining work.} Global dynamics, larger shocks, endogenous type
selection, and a policy implementing the housing upgrade remain unproved.

\textbf{Relation to Coven et al.}
Their August 2026 lifecycle model includes a minimum owner size and separate
rental and owner stocks. Their two-period illustration studies housing as an
asset without housing utility. It therefore does not supply the dated
housing-utility theorem sought here. The proposed common-stock model is a
separate illustrative exercise.\footnote{Joshua Coven, Sebastian Golder,
Arpit Gupta and Abdoulaye Ndiaye, \emph{Property Taxes and Housing Allocation
Under Financial Constraints}, August 1, 2026, pp. 15, 20--23.
\url{https://abdouecon.github.io/research/papers/Property_Tax.pdf}}

\end{document}
